% Rates of Convergence of Actor-Critic Methods: the five write-ups in docs/ as one book.
% The chapters are generated from docs/<name>.tex by make_book.py; edit those sources.
% Build with docs/build.sh (no arguments), or from docs/book:
%   python3 make_book.py && tectonic --outdir .. book.tex
\documentclass[11pt,oneside]{book}

\usepackage[T1]{fontenc}
\usepackage{lmodern}
\usepackage{amsmath,amssymb,amsthm}
\usepackage{mathtools}
\usepackage{booktabs}
\usepackage{graphicx}
\usepackage[nottoc]{tocbibind}
\usepackage[margin=1.1in]{geometry}
\usepackage[colorlinks=true,linkcolor=blue,citecolor=blue,urlcolor=blue,
  bookmarksnumbered]{hyperref}
\emergencystretch=3em
\sloppy

% Numbering: chapter.section.n, so Theorem 4.7 of a chapter's note is Theorem c.4.7 here.
\newtheorem{theorem}{Theorem}[section]
\newtheorem{lemma}[theorem]{Lemma}
\newtheorem{proposition}[theorem]{Proposition}
\newtheorem{corollary}[theorem]{Corollary}
\theoremstyle{definition}
\newtheorem{definition}[theorem]{Definition}
\newtheorem{assumption}[theorem]{Assumption}
\newtheorem{fact}[theorem]{Cited Fact}
\newtheorem{observation}[theorem]{Observation}
\theoremstyle{remark}
\newtheorem{remark}[theorem]{Remark}
\numberwithin{equation}{chapter}

% The union of the chapters' macros.
\newcommand{\R}{\mathbb{R}}
\newcommand{\E}{\mathbb{E}}
\newcommand{\F}{\mathcal{F}}
\newcommand{\X}{\mathcal{S}}
\newcommand{\U}{\mathcal{A}}
\newcommand{\Y}{\mathcal{Y}}
\newcommand{\Lop}{\mathcal{L}}
\newcommand{\IV}{\mathrm{IV}}
\newcommand{\one}{\mathbf{1}}
\newcommand{\inner}[2]{\langle #1, #2\rangle}
\DeclareMathOperator{\tr}{tr}
\DeclareMathOperator{\spn}{span}
\DeclareMathOperator{\adj}{adj}
\DeclareMathOperator{\sgn}{sgn}
\newcommand{\lean}[1]{\texttt{\detokenize{#1}}}
\newcommand{\leanref}[1]{\par{\raggedleft\small Lean:~\lean{#1}\par}}

% The chapter's opening summary (the abstract of its note), and a pointer to its source.
\newenvironment{chapterabstract}{\begin{quote}\small}{\end{quote}\medskip}
\newcommand{\chaptersource}[1]{\noindent{\small\itshape Source: \texttt{docs/#1.tex}
  (Markdown: \texttt{docs/#1.md}).}\par\medskip}

\title{Rates of Convergence of Actor-Critic Methods\\[6pt]
  \large Asymptotic covariances, critics and eligibility traces,\\
  with a formalization in Lean~4}
\author{Vijaymohan Konda}
\date{November 2026}

\begin{document}
\pagenumbering{Alph}
\maketitle
\frontmatter

\chapter{Preface}

These notes study how fast actor-critic methods converge on finite discounted Markov
decision processes, and how the rate depends on the critic: its features, and the trace
parameter $\lambda$ of TD($\lambda$). The figure of merit is the loss in expected reward,
$\E[J(\theta^*) - J(\theta_k)]$. To first order in the step size it is the asymptotic covariance of the
actor's parameter, weighted by the Hessian of $J$. The comparison with actor-only methods
(REINFORCE) is the thread through the book.

\paragraph{Structure.}
\begin{itemize}
\item Chapter~\ref{ch:pg} proves the policy gradient theorem for finite MDPs and the regularity of
  $\nabla J$ used later.
\item Chapter~\ref{ch:td} studies TD($\lambda$) for policy evaluation: the bias of its limit and its
  intrinsic variance, with closed forms for finite chains.
\item Chapters~\ref{ch:cov} and~\ref{ch:tts} give the covariance theory. For linear stochastic
  approximation with Markov noise, on one and on two time scales, $\beta_k^{-1}\E[\theta_k\theta_k^\top]$
  converges to the solution of a Lyapunov equation.
\item Chapter~\ref{ch:ac} puts these together:
  \begin{itemize}
  \item it linearizes the actual iterations, with Markov and $\theta$-dependent sampling;
  \item it relates the covariance to the loss;
  \item it compares critics. A TD(1) critic matches REINFORCE. A critic that represents the value
    function is unbiased for every $\lambda$. For the tabular critic, shorter traces give smaller
    noise, so TD(0) is the best tabular critic and TD(1) the worst.
  \end{itemize}
\end{itemize}
Chapter~\ref{ch:ac} uses $\nabla J$ from Chapter~\ref{ch:pg} and the covariance theorems of
Chapters~\ref{ch:cov} and~\ref{ch:tts}. The trace-driven noise of Chapter~\ref{ch:cov} uses the
closed-form covariance of Chapter~\ref{ch:td}.

\paragraph{Formalization.} Most results are machine-checked in Lean~4, in the
\texttt{RLTheory/} library of the repository \texttt{github.com/vmkonda/rl-theory-in-lean}, a
fork of \texttt{ShangtongZhang/rl-theory-in-lean}. A formalized statement is followed by a
\emph{Lean} note naming its theorem. The hypotheses of the Lean theorem are the assumptions of
the result. A statement without such a note is proved on paper only. Results marked
\emph{cited facts} are used without proof. Results marked \emph{observations} rest on
computation, with code in \texttt{experiments/}. Appendix~\ref{app:lean} explains how to check a
theorem and lists the Lean modules by chapter. The last section of Chapter~\ref{ch:ac} lists what
is and is not formalized. The main open item is global dissipativity, which a local theory would
replace by a localization argument.

\paragraph{Numbering.} Results are numbered chapter.section.n. Each chapter also exists as a
standalone note in \texttt{docs/}, numbered section.n, which is the numbering used in the Lean
docstrings. For example, Theorem~4.7 of \texttt{docs/actor-critic-rates.md} is
Theorem~\ref{ac:thm:mono} here. References inside a chapter to ``the companion note'' or to
another note mean the corresponding chapter of this book.

\tableofcontents

\mainmatter
\input{build/ch-pg}
\input{build/ch-td}
\input{build/ch-cov}
\input{build/ch-tts}
\input{build/ch-ac}

\appendix
\chapter{Checking the Formalization}\label{app:lean}

\paragraph{Building.} From the root of the repository, \texttt{lake exe cache get} fetches the
compiled Mathlib, and \texttt{lake build} checks every file. The build fails on any error, and a
\texttt{sorry} appears as a warning. The library builds with none.

\paragraph{Axioms.} A theorem rests only on the standard axioms when
\begin{verbatim}
import RLTheory.Basic
#print axioms StochasticApproximation.ScoreChain.restart_actor_loss_mono
\end{verbatim}
in a scratch file, run with \texttt{lake env lean <file>}, prints
\texttt{[propext, Classical.choice, Quot.sound]}. Every theorem named in this book does.

\paragraph{Modules by chapter.} All paths are under \texttt{RLTheory/}.

\medskip
\begin{tabular}{@{}lp{0.58\textwidth}@{}}
\toprule
Chapter & Modules \\
\midrule
\ref{ch:pg} & \texttt{MarkovDecisionProcess/PolicyFamily}, \texttt{PolicyGradient},
  \texttt{PolicyGradientLipschitz}; \texttt{StochasticApproximation/GradientAscent} \\
\ref{ch:td} & none: proved on paper, with numerics in \texttt{experiments/td\_lambda.py} \\
\ref{ch:cov} & \texttt{StochasticApproximation/CovarianceRecursion}, \texttt{MarkovNoiseCovariance},
  \texttt{PoissonNoise}, \texttt{TraceMoments}, \texttt{MultiplicativeNoiseCovariance} \\
\ref{ch:tts} & \texttt{StochasticApproximation/TwoTimeScaleRecursion}, \texttt{TwoTimeScaleCovariance},
  \texttt{TwoTimeScaleMultiplicative}, \texttt{TwoTimeScaleTrace} \\
\ref{ch:ac} & \texttt{StochasticApproximation/Linearization}, \texttt{TwoTimeScaleNonlinear},
  \texttt{TwoTimeScaleMoments}, \texttt{TwoTimeScaleMarkov}, \texttt{ParamMarkovNoise},
  \texttt{DoeblinPoisson}, \texttt{DoeblinCovariance}, \texttt{TwoTimeScaleParam},
  \texttt{RateComparison}, \texttt{OrthogonalNoise}, \texttt{LookbackMonotone};
  \texttt{MarkovDecisionProcess/CriticTheorems}, \texttt{RestartProcess}, \texttt{LambdaMonotone} \\
\bottomrule
\end{tabular}

\backmatter
\input{build/bibliography}

\end{document}
